\section{Deklaration des KI-Einsatzes}\label{app:ki-deklaration}

Im Rahmen dieser Arbeit wurden Claude Code von Anthropic und GPT-5.6 Sol
von OpenAI eingesetzt. Claude Code wurde als Agent für Agentic Coding
verwendet. Der Agent erhielt konkrete Arbeitsaufträge und konnte Dateien
im Repository erstellen oder ändern, Befehle ausführen und Tests starten.
Auf diese Weise entstand der grösste Teil der Dateien im Repository. Dazu
gehören der Datengenerator, das Dataset Tool, die Trainingspipeline, die
mobilen Apps sowie Skripte, Tests und Teile der projektbegleitenden
Dokumentation.

Die Arbeit mit Claude Code erfolgte in eng begleiteten Sitzungen. Der Verfasser
definierte jeweils die Anforderungen, Ziele und Randbedingungen. Danach
prüfte er die Änderungen im Code, führte Tests aus und kontrollierte die
erzeugten Ergebnisse. Jeder Commit wurde durch den Verfasser in einem
Review geprüft. Fehler, unklare Annahmen oder Abweichungen führten zu
einem neuen und genauer formulierten Arbeitsauftrag. Alle Entscheide zur
Architektur, zu den Daten, zum Training, zur Bedienung der Apps und zum
weiteren Vorgehen wurden durch den Verfasser getroffen.

Die vorliegende Transferarbeit wurde vom Verfasser selbst geschrieben.
Für die sprachliche Überarbeitung einzelner Textstellen wurde GPT-5.6 Sol
eingesetzt. Das Modell unterstützte bei der Korrektur von Grammatik und
Satzstellung, bei der stilistischen Verbesserung sowie bei der Prüfung
auf Richtigkeit. Alle vorgeschlagenen Änderungen wurden vom Verfasser
geprüft und bei Bedarf angepasst.

KI-Ausgaben werden in dieser Arbeit nirgends als Quelle zitiert. Aussagen,
die einen Beleg benötigen, verweisen auf die im Literaturverzeichnis
aufgeführten Quellen. Die Verantwortung für den Inhalt, die Entscheide,
die Ergebnisse und die Schlussfolgerungen liegt beim Verfasser.

\begin{table}[H]
\centering\small
\begin{tabularx}{\textwidth}{@{}>{\raggedright\arraybackslash}p{3.0cm}YY@{}}
\toprule
\textbf{Bereich} & \textbf{Art des KI-Einsatzes} & \textbf{Beitrag des Verfassers} \\
\midrule
Software und Repository &
Claude Code erstellte und bearbeitete nach konkreten Arbeitsaufträgen Code,
Skripte, Tests, Workflows und technische Dokumentation. &
Definition der Anforderungen und der Architektur. Vorgabe der Schnittstellen
und des gewünschten Verhaltens. Prüfung und Freigabe jedes Commits. \\
\addlinespace
Daten, Training und Auswertung &
Claude Code unterstützte bei der Umsetzung des Datengenerators, der
Trainingspipeline, der Auswertungsskripte und der Abbildungen. &
Erhebung und Labeling der realen Daten. Festlegung der Parameter und des
Versuchsaufbaus. Durchführung der Läufe sowie Interpretation der Ergebnisse. \\
\addlinespace
Prüfung und Qualitätssicherung &
Claude Code erstellte Tests und automatisierte Prüfungen. Der Agent
unterstützte zudem bei der Suche nach Fehlern und bei möglichen Korrekturen. &
Festlegung der Prüfkriterien. Ausführung der Tests und visuelle Kontrolle der
Ergebnisse. Beurteilung der Befunde und Entscheid über jede Korrektur. \\
\addlinespace
Schriftliche Arbeit &
GPT-5.6 Sol wurde zur Korrektur von Grammatik und Satzstellung sowie zur
Prüfung auf Richtigkeit eingesetzt. &
Verfassen des gesamten Textes. Festlegung von Aufbau, Inhalt,
Argumentation und Schlussfolgerungen. Prüfung und Übernahme der Korrekturen. \\
\bottomrule
\end{tabularx}
\caption{Deklaration des KI-Einsatzes nach Bereichen.}
\label{tab:ki-deklaration}
\end{table}
